\section{Introduction}\label{cc:sec:problem}

The incompressible Euler equations are among the earliest partial
differential equations derived in mathematical physics. Euler presented them in
1755, less than a decade after d'Alembert derived the wave equation for
the vibrating string~\cite{DAlembert1747,Euler1757}.

The goal of this work is to construct smooth initial data for the
incompressible Euler equations on the three-dimensional torus whose
solution blows up in finite time. We represent $\mathbb T^3$ by the
centered cube $[-1/2,1/2]^3$ with periodic boundary conditions
(opposite faces identified). We write these equations as
\begin{equation}\label{cc:eq:euler}
 \partial_tu+(u\cdot\nabla)u+\nabla p=0,
 \qquad \diver u=0
 \quad\hbox{on }\mathbb T^3\times[0,T_*).
\end{equation}
Local existence and uniqueness for sufficiently regular data go back to
Lichtenstein and G\"unther in the 1920s~\cite{Lichtenstein1925,Gunther1926}.
Classical theory on $\mathbb T^3$ supplies a unique maximal smooth solution
from smooth divergence-free data~\cite{EbinMarsden,BedrossianVicol}.
By the Beale--Kato--Majda criterion~\cite{BKM,GGPP}, a smooth solution on
$[0,T)$ continues smoothly through $T<\infty$ if
$\int_0^T\|\curl u(t)\|_{L^\infty}\,\dd t<\infty$.

In two dimensions the scalar vorticity is transported, so its
$L^\infty$ norm is conserved, and smooth solutions exist for all
time~\cite{Wolibner1933}. Yudovich~\cite{Yudovich1963} proved global
existence and uniqueness for bounded, possibly nonsmooth initial
vorticity. In three dimensions, Fefferman's description of the
Navier--Stokes Millennium Problem stated that the global regularity
question on $\R^3$ and on $\mathbb T^3$ was also open for the Euler
equations~\cite{Fefferman2006}.

Numerical searches for blow-up include the Taylor--Green vortex study
of Morf, Orszag and Frisch~\cite{MorfOrszagFrisch1980} and Kerr's
computations for antiparallel vortex tubes~\cite{Kerr1993}.
In a later computation, Hou and Li~\cite{HouLi2006} observed at most
double-exponential growth of the maximum vorticity past Kerr's predicted
singular time.
Luo and Hou~\cite{LuoHou2014} found numerical evidence for axisymmetric
blow-up at the wall of an axially periodic cylinder.
Childress, Gilbert and Valiant~\cite{ChildressGilbertValiant2016}
studied, analytically and numerically, eroding vortex dipoles in
axisymmetric flow without swirl, whose maximal vorticity grows like
$t^{4/3}$, and Childress and Gilbert~\cite{ChildressGilbert2018}
proposed, by formal asymptotic analysis, the hairpin geometry as a model
for finite-time blow-up.

On $\R^3$, Elgindi~\cite{ElgindiSingularity} constructed finite-time
blow-up from $C^{1,\alpha_0}$ initial data for some $\alpha_0>0$.
Elgindi, Ghoul and Masmoudi~\cite{ElgindiGhoulMasmoudi} obtained stable,
locally self-similar blow-up for finite-energy data in this regularity class.
Chen and Hou~\cite{ChenHouSmoothI,ChenHouSmoothII} claim nearly self-similar
blow-up from smooth axisymmetric data in a periodic cylinder with boundary,
using computer-assisted estimates. The theorem below concerns smooth
periodic data on $\mathbb T^3$, with bounded velocity and vorticity growth
concentrated at the origin.

\begin{theorem}\label{cc:target:unforced}
There is an odd divergence-free
$u_0\in C^\infty(\mathbb T^3;\R^3)$ whose maximal classical solution $u$
of \eqref{cc:eq:euler} is smooth and odd on
$\mathbb T^3\times[0,T_*)$ for a finite $T_*>0$, and satisfies
\begin{equation}\label{cc:uf:eq:main-breakdown}
 \int_0^{T_*}\|\nabla u(t)\|_{L^\infty}\dd t=\infty,
 \qquad
 \int_0^{T_*}\|\curl u(t)\|_{L^\infty}\dd t=\infty.
\end{equation}
Moreover,
\begin{equation}\label{cc:uf:eq:main-terminal}
 \sup_{0\le t<T_*}\|u(t)\|_{L^\infty}<\infty,
 \qquad \lim_{t\uparrow T_*}|\curl u(t,0)|=\infty.
\end{equation}
For every compact $K\subset\mathbb T^3\setminus\{0\}$ and every integer
$m\ge0$, both $u(t)$ and $\partial_tu(t)$ converge in $C^m(K)$ as
$t\uparrow T_*$.
\end{theorem}

\begin{remark}
The solution in Theorem~\ref{cc:target:unforced} may be chosen so that it
extends smoothly in space and time through $T_*$ away from the origin.
A proof is given in the auxiliary note~\cite{TerminalSmoothness}.
\end{remark}

We construct a sequence of smooth approximate solutions of the Euler
equations whose residuals tend to zero. The sequence converges in every
spatial $C^m$ norm, uniformly on compact time intervals before $T_*$, to
a solution of the Euler equations.

At each step in the sequence we add an oscillation at a higher frequency, localized
near the origin and initially small in velocity. The velocity gradient
of the preceding scales amplifies the new oscillation, whose gradient
in turn drives the growth of the next one. The growth intervals have
finite total length. This multiscale construction follows the approach of C\'ordoba and
Mart\'inez-Zoroa for the Euler equations with forcing and for forced
IPM~\cite{CMZForcedEuler,CMZForcedIPM}.
C\'ordoba, Mart\'inez-Zoroa and Zheng~\cite{CMZUnforcedEuler} used a cascade of
vorticity layers to obtain blow-up for the unforced Euler equations
on $\R^3$ from $C^{1,\alpha}$ data
smooth away from the origin. Related constructions give blow-up with forcing
for the inviscid Boussinesq equations, the variable-density Euler equations
and the hypodissipative Navier--Stokes equations~\cite{CLMZForcedBoussinesq,CLMZNonhomogeneousEuler,CMZHypodissipativeNS}.
Our companion papers proving singularity formation for forced IPM,
the forced Boussinesq equations, and the forced Euler
equations~\cite{AlpogeBuckmasterCoiculescuIPM,AlpogeBuckmasterBoussinesq,AlpogeBuckmasterForcedEuler}
use forces smooth in space and time up to the singular time. In the
present construction we use Craik--Criminale waves as building blocks
and choose the initial velocity at each step so that the limiting
solution is unforced.

Convex integration constructions of nonunique weak solutions of the Euler equations
also use oscillations at successive frequency scales. Below the Onsager
exponent, this method yields solutions with nonconstant kinetic
energy~\cite{IsettOnsager,BDLSVOnsager}. For the unforced Euler equations
in the plane, Coiculescu~\cite{CoiculescuEulerI,CoiculescuEulerII}
uses nested vortex rings at different scales. An instability propagates
from the inner rings outward, yielding distinct renormalized weak
solutions with the same initial vorticity. See also Choi~\cite{ChoiEulerNonuniqueness}.
For the three-dimensional Navier--Stokes equations, Jia and
\v{S}ver\'ak~\cite{JiaSverak2015} gave a conditional construction of
nonunique Leray--Hopf solutions based
on spectral assumptions about forward self-similar flows.
Nonuniqueness was shown for
finite-energy weak solutions by Buckmaster and Vicol~\cite{BuckmasterVicolNS},
and in the Leray--Hopf class with a common force by Albritton, Bru\'e and
Colombo~\cite{AlbrittonBrueColombo}. In a recent preprint, Buckmaster,
Coiculescu and Li~\cite{BuckmasterCoiculescuLiNS} obtain nonuniqueness
of unforced Leray--Hopf solutions.

\subsection{Waves and successive amplification}\label{cc:sec:intro-waves}
We first describe the wave construction of Craik and
Criminale~\cite{CraikCriminale}. Let $D(t)$ be a smooth matrix with
$\operatorname{tr}D=0$ and $D'+D^2$ symmetric. These conditions ensure that
$U(t,x)=D(t)x$, linear in space, solves the Euler equations on $\R^3$, with pressure
$p_D(t,x)=-x\cdot(D'+D^2)x/2$. We add to this linear background a velocity
\[
 w(t,x)=\mathbf a(t)F(k(t)\cdot x),\qquad \mathbf a\cdot k=0,
\]
where $k(t)\ne0$ is the wave vector, $\mathbf a(t)$ is the amplitude
vector, and $F$ is a smooth odd $1$-periodic profile of the scalar phase $s$.
For example, one may take $F(s)=\sin(2\pi s)$. The vector $\mathbf a$
records both size and direction: the velocity amplitude is $|\mathbf a|\|F\|_\infty$, and
$\mathbf a/|\mathbf a|$ is its unit direction when $\mathbf a\ne0$.
The unit wave vector is $k/|k|$. The length $|k|$ measures the spatial
frequency in cycles per unit length for the first phase harmonic.

The sum $u=Dx+w$ is an exact solution of the Euler equations if $k$ and $\mathbf a$ solve
\begin{equation}\label{cc:intro:wave-ode}
 k'=-D^Tk,\qquad
 \mathbf a'=-D\mathbf a+p_{\rm w} k,\qquad
 p_{\rm w}=\frac{2k\cdot D\mathbf a}{|k|^2}.
\end{equation}
Its pressure is $p=p_D-p_{\rm w} P(k\cdot x)$, where $P'=F$.
Indeed, $\mathbf a\cdot k=0$ gives $(w\cdot\nabla)w=0$, the first
ordinary differential equation transports the phase by $Dx$, and the
second cancels the remaining acceleration with the gradient of the
wave pressure $-p_{\rm w}P(k\cdot x)$.
The displayed value of $p_{\rm w}$ preserves $\mathbf a\cdot k=0$.
We call \eqref{cc:intro:wave-ode} the Craik--Criminale equations for the
wave vector and the amplitude vector. Normalizing $F'(0)=1$ gives
$\nabla w(t,0)=\mathbf a\otimes k$.
Related short-wave instabilities of planar vortices were studied by
Pierrehumbert~\cite{Pierrehumbert1986} and Bayly~\cite{Bayly1986}.

To describe amplification of a new wave, write $k_p$ and $\mathbf a_p$
for the nonzero wave vector and amplitude vector of the preceding wave,
and set
\[
 E_1=\frac{k_p}{|k_p|},\qquad
 E_2=\frac{\mathbf a_p}{|\mathbf a_p|},\qquad E_3=E_1\times E_2,
 \qquad \sigma=|\mathbf a_p|\,|k_p|.
\]
$E_1$ is the unit wave vector and $E_2$ the unit polarization vector.
The preceding wave contributes the velocity gradient
$\mathbf a_p\otimes k_p=\sigma E_2\otimes E_1$ at the origin.
Here $(E_2\otimes E_1)v=(E_1\cdot v)E_2$: the corresponding linear
velocity is $\sigma(E_1\cdot x)E_2$, a shear in direction $E_2$
varying in direction $E_1$. The new wave therefore sees
$D=B+\sigma E_2\otimes E_1$, where $B$ is the older background gradient.
To model the wave equations at the origin, take $B$ with just two
nonzero entries in this moving frame, both constant,
\[
 E_1\cdot BE_2=\nu>0,\qquad E_3\cdot BE_2=\eta>0,
 \qquad \varepsilon=\eta/\nu\ll1.
\]
Start the new wave at time $t=0$ with $k=\lambda E_3$, where $\lambda>0$
is its initial frequency, and write $a_i=\mathbf a\cdot E_i$. The
preceding wave evolves by \eqref{cc:intro:wave-ode} with $D=B$.
Differentiating its normalized vectors gives, in this case,
$E_1'=-\nu E_2$, $E_2'=\nu E_1-\eta E_3$ and $E_3'=\eta E_2$.
Combining these frame equations with \eqref{cc:intro:wave-ode} for the
new wave gives, at $t=0$,
\[
 \dot a_1=-2\nu a_2,\qquad
 \dot a_2=-(\sigma-\nu)a_1.
\]
Freezing these coefficients gives a growing mode with rate
$\sqrt{2\nu(\sigma-\nu)}$ when $\sigma>\nu$. We choose the initial
polarization along $E_2-\sqrt{2\nu/\sigma}\,E_1$, the growing direction
to leading order for $\sigma\gg\nu$.

The coefficients evolve as the wave vector turns. Let $\varphi$
be the angle of $k$ from $E_3$ toward $E_1$:
\[
 \tan\varphi=\frac{k\cdot E_1}{k\cdot E_3},
\]
so $\varphi=0$ at $t=0$. The coupling $\eta$ tilts $k$ toward $-E_2$,
and the shear carries this tilt into $E_1$, so $\varphi$ increases;
while $\eta t$ is small, $\tan\varphi$ grows like $(\sigma+\nu)\eta t^2$.
In the strong-shear limit $\sigma/\nu\to\infty$, the wave pressure
multiplies the driving coefficient $2\nu$ in the equation for $a_1$ by
$\cos^2\varphi\,(1-\varepsilon\tan\varphi)$. This factor changes sign
at $\tan\varphi=1/\varepsilon$. The reference amplitude reaches its peak
after this sign change, as shown in Figure~\ref{cc:fig:wave-turning}.
For the perturbed wave with the chosen polarization,
Proposition~\ref{cc:prop:finite-return} gives a \emph{gain} of at least
a fixed multiple of $\exp(\varepsilon^{-1/2})$ on its prescribed interval.
The gain is the ratio of the maximum of $|\mathbf a|$ on that interval
to its initial value.

\begin{figure}[tbp]
\centering
\includegraphics[width=\textwidth]{../../output/figures/wave-turning/wave-turning-paper.pdf}
\caption{Turning and growth in the case of two nonzero entries, in the
limit of a strong shear.
Left: the angle $\varphi$ of the wave vector from $E_3$ toward $E_1$.
Right: amplitude as a fraction of its peak. The dashed line marks
the amplitude peak, by which time $\varphi$ is close to $\pi/2$.}
\label{cc:fig:wave-turning}
\end{figure}

Size alone does not determine the successor wave. The new wave must
supply both a large shear and a suitable frame for the next step of
the iteration. In its frame (the subscript $+$ denotes the frame of
the newly added wave),
\[
 E_{1,+}=k/|k|,\qquad E_{2,+}=\mathbf a/|\mathbf a|,\qquad
 E_{3,+}=E_{1,+}\times E_{2,+},
\]
its shear is $\sigma_+E_{2,+}\otimes E_{1,+}$, with
$\sigma_+=|\mathbf a|\,|k|$. The older gradient $D$, expressed in this
new frame, must again have positive couplings
\[
 \nu_+=E_{1,+}\cdot D E_{2,+},\qquad
 \eta_+=E_{3,+}\cdot D E_{2,+}.
\]
These play the roles of $\nu$ and $\eta$ in the next amplification
step. We prescribe a smaller ratio $\eta_+/\nu_+=\varepsilon_+<\varepsilon$
to give the next wave enough amplification to keep the error from its
time cutoff small.
The turning of the wave vector supplies this orientation. In the
reference calculation $\nu_+$ starts at $\varepsilon\nu$, and the
large terms in the numerator for $\eta_+$ cancel. For the perturbed wave,
Lemma~\ref{cc:lem:finite-comparison} shows that at late times $\nu_+$
is close to $\nu$ and $\eta_+/\nu_+$ is comparable to $\cos\varphi$.
Under the hypotheses of
Proposition~\ref{cc:prop:finite-return}, the ratio reaches the prescribed
value after the reference amplitude peaks, while the new shear remains
large. We start the successor at that crossing.
Thus the Craik--Criminale equations produce both amplification and the
orientation needed to repeat the construction.

\Needspace{10\baselineskip}
The periodic construction starts from the stationary shear flow
\[
 U_0(x)=\sin(2\pi x_2)(e_1+s e_3),
\]
where $e_i$ are the coordinate vectors and $s>0$ is small. It solves
the Euler equations with zero pressure.
In the initial coordinate frame its central gradient has
$\nu=2\pi$ and $\eta=2\pi s$. Its explicit flow supplies the first wave;
subsequent waves use the successive gradients they create.

Craik--Criminale waves on $\R^3$ provide a local model for amplification.
For the periodic construction, we localize each wave by a smooth cutoff
supported in the interior of a fundamental cell in Lagrangian
coordinates. We extend the localized field periodically and transport its
phase and cutoff by the preceding periodic velocity, as in geometric optics
\cite{LifschitzHameiri1991,FriedlanderVishik1991,Lax1957,ChoquetBruhat1969}.
Since the cutoff
vanishes near the cell boundary, the wave vector need not lie
in $\mathbb Z^3$. Localization and the spatial variation of the background
introduce errors, which the corrections below make small at each step
and remove in the limit.

\subsection{Controlling corrections in an amplifying flow}
To correct a velocity in an amplifying flow, we use the freedom to choose
its initial value. To avoid exciting growing modes, we prescribe the final
particle positions and choose the initial velocity to reach them.

Let $u,p$ solve the Euler equations on $[0,T]$. The flow $X(t,a)$ of
$u$ gives the position at time $t$ of the particle that starts at $a$.
It satisfies
\[
 \partial_tX(t,a)=u(t,X(t,a)),\qquad X(0,a)=a.
\]
For a divergence-free velocity correction $b$, let $X_\delta$ be the
flow of $u+\delta b$ with the same initial positions, so that
\[
 \partial_tX_\delta(t,a)
       =(u+\delta b)(t,X_\delta(t,a)),\qquad X_\delta(0,a)=a.
\]
The first-order change $\xi$ in the particle trajectories is defined by
\[
 X_\delta(t,a)=X(t,a)+\delta\xi(t,a)+O(\delta^2).
\]
Differentiating the flow equation at $\delta=0$ gives
\begin{equation}\label{cc:intro:first-variation}
 \partial_t\xi(t,a)=\nabla u(t,X(t,a))\xi(t,a)+b(t,X(t,a)),
 \qquad \xi(0,a)=0.
\end{equation}
The initial condition follows from the common starting positions.
We impose $\xi(T,a)=0$ to select $b(0)=\partial_t\xi(0)$ so that the final
particle positions agree to first order.

For a prescribed momentum source $f$, the inhomogeneous problem is to
find $b,q$ satisfying the linearized Euler equations
\begin{equation}\label{cc:intro:linearized-momentum}
 \partial_t b+(u\cdot\nabla)b+(b\cdot\nabla)u+\nabla q=f,
 \qquad \diver b=0.
\end{equation}
Here $q$ is the pressure variation, and the displacement induced by $b$
must satisfy $\xi(0)=\xi(T)=0$. Differentiating
\eqref{cc:intro:first-variation} in time and using
\eqref{cc:intro:linearized-momentum} gives
\begin{equation}\label{cc:intro:jacobi}
 \begin{gathered}
 \xi_{tt}+H\xi+(\nabla q)\circ X=f\circ X,
 \qquad H=(D_x^2p)\circ X,\\
 \diver_x(\xi\circ X^{-1})=0,\qquad \xi(0)=\xi(T)=0.
 \end{gathered}
\end{equation}
The homogeneous equation is the Jacobi equation along the geodesic
associated with this solution in Arnold's formulation~\cite{Arnold1966}.
The constraint $\diver_x(\xi\circ X^{-1})=0$ expresses volume
preservation to first order.

For a solution of \eqref{cc:intro:jacobi}, testing with $-\xi$ and
integrating by parts gives the identity
\begin{equation}\label{cc:intro:Q-identity}
 \underbrace{\int_0^T\!\int_{\mathbb T^3}
   \bigl(|\xi_t|^2-\xi\cdot H\xi\bigr)\,\dd a\,\dd t}_{\mathcal Q(\xi)}
 =-\int_0^T\!\int_{\mathbb T^3}(f\circ X)\cdot\xi\,\dd a\,\dd t.
\end{equation}
Volume preservation and the divergence constraint remove the pressure
variation. The quadratic form $\mathcal Q$ is the fixed-endpoint second variation of the action.
If $H\le\Id$ and $T\le1$, then
${|\xi(t,a)|^2\le(T-t)\int_t^T|\xi_\tau(\tau,a)|^2\,\dd\tau}$ gives
\begin{equation}\label{cc:intro:coercivity}
 \mathcal Q(\xi)\ge\Bigl(1-\frac{T^2}{2}\Bigr)\|\xi_t\|_{L^2_{t,a}}^2
 \ge\frac12\|\xi_t\|_{L^2_{t,a}}^2 .
\end{equation}
The term $-\xi\cdot H\xi$ lowers the integrand where $\xi\cdot H\xi>0$
and raises it where $\xi\cdot H\xi<0$. Thus an upper bound for $H$ gives
the lower bound in \eqref{cc:intro:coercivity}.
For $f=0$, \eqref{cc:intro:Q-identity} gives $\mathcal Q(\xi)=0$. Then
\eqref{cc:intro:coercivity} forces $\xi_t=0$, and $\xi(0)=0$ gives $\xi=0$.
Equation~\eqref{cc:intro:first-variation} gives $b=0$, proving that the linearized
endpoint map $b(0)\mapsto\xi(T)$ has trivial kernel.
Upper bounds on the pressure Hessian also enter Brenier's fixed-endpoint
action-minimization criterion~\cite{Brenier1989,Brenier1999}.

For an approximate parent $U$ with residual $R_U$,
Subsection~\ref{cc:uf:sec:parent-step} uses
$H^\#=(D_x^2p_U-\nabla R_U)\circ X$ in place of $H$, under the same
bounds $H^\#\le\Id$ and $T\le1$. Subsection~\ref{cc:uf:sec:mean} uses
the resulting coercivity to construct a bounded inverse for the forced
Jacobi problem, and
Appendix~\ref{cc:app:mean} establishes the quantitative H\"older bounds
of Theorem~\ref{cc:uf:mean:inverse}.

In the affine wave of Subsection~\ref{cc:sec:intro-waves}, the Hessian of
$-p_{\rm w}P(k\cdot x)$ is $-p_{\rm w}F'(k\cdot x)k\otimes k$.
We seek an odd periodic profile
$F_h$ satisfying
\begin{equation}\label{cc:intro:profile-slopes}
 F_h'(0)=1,\qquad -Ch\le F_h'(s)\le1 .
\end{equation}
The first condition preserves the wave gradient $\mathbf a\otimes k$
at the origin, and hence the amplification described in Subsection~\ref{cc:sec:intro-waves}.
The lower bound on $F_h'$ controls the largest eigenvalue of
$D_x^2[-p_{\rm w}P_h(k\cdot x)]$, where $P_h'=F_h$. On an interval
with $p_{\rm w}>0$, its positive part satisfies
\begin{equation}\label{cc:intro:profile-curvature}
 p_{\rm w}|k|^2\bigl(-F_h'(k\cdot x)\bigr)_+
 \le Ch\,p_{\rm w}|k|^2,
\end{equation}
where $(r)_+=\max\{r,0\}$. For fixed wave coefficients, taking $h$ small
makes this positive part small at every point. In the construction $h$
is a negative power of the frequency of the wave. As $h$ shrinks, the
wave produces thin layers of strong shear while its velocity amplitude
is $O(h|\mathbf a|)$. Its gradient at the origin remains
$\mathbf a\otimes k$, preserving the local shear that amplifies the
next wave. The pressure Hessian added at each step is bounded above
by a small multiple of the identity. Section~\ref{cc:uf:sec:tower}
sums these bounds and controls $\|\nabla R_U\|_\infty$ to obtain
$H^\#\le\Id$ as a quadratic form.

For $F(s)=\sin(2\pi s)$, the ratio $F'(1/2)/F'(0)$ is $-1$.
The sine profile with slope one at the origin therefore violates the
lower bound in \eqref{cc:intro:profile-slopes} for small $h$.
Subsection~\ref{cc:sec:profile} constructs $F_h$ with a narrow positive
bump in $F_h'$ and a shallow negative part, as shown in
Figure~\ref{cc:fig:pressure-profile}.

\begin{figure}[!hbp]
\centering
\includegraphics[width=.75\textwidth]{../../output/figures/pressure-and-timeline/pressure-profile.pdf}
\caption{The sine slope divided by its value at the origin (dashed)
and the slope $F_h'$ of the modified profile (solid). Both curves equal
one at the origin.}
\label{cc:fig:pressure-profile}
\end{figure}

\subsection{Correcting a localized wave}
We smoothly truncate the velocity of each step in frequency at the end of the step,
to simplify the bookkeeping of derivative estimates, as in the mollification
step of convex integration~\cite[Section~2.4]{BDLSVOnsager}.
This introduces a small error, carried to the next step;
Subsection~\ref{cc:uf:trunc:sec:truncation} gives the estimates.
Write $U$ for the truncated velocity of the preceding step and $R_U$ for
its residual in the Euler equations.

We add the localized wave described in Subsection~\ref{cc:sec:intro-waves},
using $F_h$ and a divergence correction. Let $X$ be the flow of $U$ from
time zero, and fix a nonzero vector $k_*$. The phase $\theta$ of the
localized wave is defined by
\[
 \theta(t,X(t,a))=k_*\cdot a
\]
for Lagrangian coordinates $a$ in the fundamental cell. It is the plane
phase $k_*\cdot x$ at time zero, transported by $U$. Its gradient
$\kappa=\nabla_x\theta$ is the wave vector of the localized wave; for a
linear flow, $\theta=\kappa\cdot x$, as in
Subsection~\ref{cc:sec:intro-waves}. For a wave after the first, $k_*$
is chosen so that $\kappa=\lambda E_3$ at the origin at its activation
time, as for the new wave of Subsection~\ref{cc:sec:intro-waves}.
The leading velocity is $A_1(t,x,\theta(t,x))$. Writing the profile as
$A_1(t,x,s)$ with $s$ independent lets us distinguish derivatives of
the coefficients from the large term $\kappa\partial_s$ that arises
when differentiating through the phase.
A wave after the first is multiplied
by a cutoff $\widetilde\chi(t)$ in time that vanishes before its
amplification interval; for the first wave $\widetilde\chi=1$.
The derivative of the cutoff leaves an error of the size of the wave at
the start of that interval, also carried to the next step.

The velocity added to $U$ is an expansion
\[
 \sum_{p\ge1}Z_p,\qquad Z_p=V_p+C_{p-1}+B_p,
\]
where $V_p(t,x,s)$ is a profile perpendicular to $\kappa$ with zero
phase average, $C_p$ is the divergence correction of $V_p$, and
$B_p(t,x)$ is a mean, independent of the phase; the profiles are
evaluated at $s=\theta$. The leading term is the wave,
$Z_1=V_1=\widetilde\chi A_1$, with $C_0=0$ and $B_1=0$. We call $p$ the
\emph{level} of $Z_p$: up to a fixed power of a number that bounds the
parent, the bound for $Z_p$ is that of the wave times
$\lambda^{-2(p-1)/3}$, where $\lambda$ is the frequency of the wave.
The error left by the wave and its divergence correction $C_1$ is a
profile $H_2(t,x,s)$ of level two.
We split $H_2$ into its phase average $\av{H_2}$ and its
oscillatory part $H_2-\av{H_2}$, as in the expansions of
Cheverry~\cite[Section~4.1]{Cheverry2006}.

The mean correction $B_2$ cancels the phase average of $H_2$ together
with the parent residual. With $D_t=\partial_t+U\cdot\nabla_x$, the
linearized Euler equation \eqref{cc:intro:linearized-momentum} gives
$D_t B_2+(\nabla U)B_2+\nabla q=-\av{H_2}-R_U$ with
$\diver B_2=0$. The Jacobi inverse selects $B_2(0)$ so that the
correction leaves the truncated parent's final particle positions fixed
to first order.

With $B_2$ determined, the oscillatory correction $V_2$ and its pressure
$\pi_2$ solve
\begin{equation}\label{cc:intro:transverse-correction}
 \begin{split}
 D_tV_2+(\nabla U)V_2+\kappa\partial_s\pi_2
   &=-(H_2-\av{H_2})-(B_2\cdot\kappa)\partial_s(\widetilde\chi A_1),\\
 \kappa\cdot V_2&=0,\qquad \av{V_2}=0.
 \end{split}
\end{equation}
The mean $B_2$ transports the leading wave's fast phase. This term has
zero phase average but is of the same order as the oscillatory part
$H_2-\av{H_2}$, so we solve for the mean $B_2$ before the oscillatory
correction $V_2$. At fixed phase along each trajectory,
\eqref{cc:intro:transverse-correction} is a forced Craik--Criminale
amplitude equation. Every wave uses
zero initial data for this correction. The curl construction of
Subsection~\ref{cc:sec:material} gives the divergence correction $C_2$
of $V_2$, so that $V_2+C_2$ is exactly divergence free. The term $C_2$
belongs to level three, where we solve for the mean and oscillatory
corrections again.

We stop the expansion at a level $P$, its depth, keeping the divergence
correction $C_P$ of $V_P$. The remainder then has a small factor
$\lambda^{-2P/3}$, once $\lambda$ is larger
than a fixed power of a number that bounds the parent
(Theorem~\ref{cc:uf:hier:thm:depth}). The truncated parent, localized wave
and linear corrections give an approximate solution $V$ of the Euler
equations. We truncate $V$ in frequency. Three choices make the three
errors of the step smaller than a fixed negative power of the frequency
of the next wave. The depth $P$ controls the remainder of the
expansion. The frequency cutoff controls the truncation error. The
ratio $\varepsilon$ of Subsection~\ref{cc:sec:intro-waves} sets the
gain, which controls the error of the cutoff in time. This cutoff acts
while the wave's amplitude is at most a fixed multiple of
$\exp(-\varepsilon^{-1/2})$ times its peak. These three errors make up
the residual of the truncated velocity, which is the parent of the next step
(Proposition~\ref{cc:uf:step:thm:step}).
The pressure and residual-gradient estimates give coercivity for
the mean equation of the next step. The leading wave and oscillatory
corrections stay in the transported cutoff region. The mean corrections
and the truncation difference may extend over the torus. Figure~\ref{cc:fig:correction-flow} summarizes the step.

\begin{figure}[tbp]
\centering
\begin{tikzpicture}[
  font=\small,
  block/.style={draw=black!65, rounded corners=2pt, align=center,
                inner xsep=5pt, inner ysep=4pt},
  correction/.style={block, text width=37mm, minimum height=8mm},
  line/.style={-{Stealth[length=1.7mm]}, semithick, draw=black!75}
]
\node[block] (parent) at (0,0) {Parent with residual $R_U$};
\node[block] (wave) at (0,-1.2)
  {Add localized wave and divergence-free corrector (\S\ref{cc:sec:material})\\
   \emph{Amplify vorticity at the origin}};
\node[inner sep=0pt] (linear-title) at (0,-2.55)
  {Linear corrections (\S\ref{cc:uf:sec:hierarchy})};
\node[correction] (mean) at (-4.65,-3.20)
  {Mean (\S\ref{cc:uf:sec:mean})};
\node[correction] (transverse) at (0,-3.20)
  {Oscillatory};
\node[correction] (curl) at (4.65,-3.20)
  {Divergence-free};
\node[draw=black!40, dashed, rounded corners=3pt,
      fit=(linear-title)(mean)(transverse)(curl), inner sep=4pt] (linear) {};
\node[block] (completion) at (0,-4.50)
  {Truncate in frequency (\S\ref{cc:uf:trunc:sec:truncation})};
\node[block] (exact) at (0,-5.30)
  {Next parent (\S\ref{cc:uf:trunc:sec:truncation})};
\draw[line] (parent) -- (wave);
\draw[line] (wave) -- (linear);
\draw[line] (mean) -- (transverse);
\draw[line] (transverse) -- (curl);
\draw[line] (linear) -- node[right, xshift=5pt]
  {Small remaining error} (completion);
\draw[line] (completion) -- (exact);
\end{tikzpicture}
\caption{One step of the construction.}
\label{cc:fig:correction-flow}
\end{figure}

\subsection{Comparison with blow-up constructions with forcing}
The construction of Alp\"oge and Buckmaster
\cite[Section~1.2]{AlpogeBuckmasterBoussinesq} also uses coupled wave
amplitudes to amplify a gradient and prepare the next wave. There the
temperature gradient grows together with a vorticity amplitude. The
accompanying shear would obstruct the next growth stage, so an external
force rotates the temperature gradient and wave direction together.
This preserves the coupling that drives the temperature amplitude while
changing the horizontal component of the wave vector, and hence the sign
of vorticity production. A scalar shooting parameter selects a rotation
pulse that returns the vorticity amplitude and that horizontal component
to zero, retaining the temperature gain. Keeping the wave vector vertical
then holds both amplitudes fixed while the next wave grows.

The construction of Alp\"oge and Buckmaster
\cite[Section~1.2]{AlpogeBuckmasterForcedEuler} for the Euler equations
with forcing builds on this mechanism for axisymmetric flow with swirl. Circulation replaces temperature and
couples to azimuthal vorticity. The angular control returns the principal
azimuthal-vorticity amplitude at the material center to zero while
retaining the circulation-gradient gain. The radial dependence of the
elliptic operator and of this coupling requires amplitudes that vary
across the transported support, with transport by the preceding velocity,
including its corrections.

In the present unforced construction, the amplified shear itself drives the successor. Its orientation
is obtained from the Craik--Criminale wave evolution; the pressure coefficient
is determined by incompressibility and is not a rotation control.
The next wave is activated at the start of its amplification interval,
with its amplitude in the growing direction (Subsection~\ref{cc:uf:sec:later-wave}).

Localization presents a further difference. In the construction for
the Boussinesq equations the profiles are linear near zero, so each layer's velocity and temperature are exactly linear
on a region containing the later layers. The streamfunction corrections vanish there, preserving
this region. In both forced constructions, the remaining errors are
absorbed by a force smooth through the singular time. To solve the unforced Euler equations,
we must remove those errors. The pressure and velocity
corrections need not preserve a neighborhood on which the velocity is linear;
the next wave uses the central gradient and estimates on its spatial variation.

Section~\ref{cc:sec:waves} proves amplification and its stability.
Sections~\ref{cc:sec:localized}--\ref{cc:sec:exact} localize the waves
and construct the linear corrections.
Sections~\ref{cc:sec:iteration} and~\ref{cc:uf:sec:tower} start and
iterate the construction, and Section~\ref{cc:sec:limit} passes to the limit.

\subsection*{Conventions}
A time supremum is indicated by $C_t$ or by the time interval. The
Leray projection is
$\PP=I-\nabla\Delta^{-1}\diver$, where $\Delta^{-1}$ acts on zero-mean
scalars. The periodic velocities
constructed below are odd and hence have zero spatial mean.
Pressure is normalized to have zero spatial mean. Spatial differences use the periodic
metric, and flows are lifted so that their displacements are periodic.
We use $(\nabla U)_{ij}=\partial_jU_i$, and for a matrix field $A$,
$(\diver A)_i=\sum_j\partial_jA_{ij}$.

